\section{Recursos}

Para el desarrollo del presente proyecto de tesis, se contará con una combinación de herramientas de software, plataformas de simulación en tiempo real y equipos físicos. A continuación, se describen los recursos principales requeridos y su función específica dentro del proyecto:

\begin{itemize}
    \item \textbf{MATLAB/Simulink:} Se utilizará para la simulación inicial del convertidor bidireccional en configuración interleaved. Permitirá validar el comportamiento teórico del sistema, afinar parámetros de diseño y modelar las estrategias de control antes de su implementación en la plataforma HIL.

    \item \textbf{Plataforma OPAL-RT:} Se empleará como entorno de simulación en tiempo real para validar el comportamiento del convertidor bajo condiciones similares a las reales, permitiendo probar la respuesta del sistema frente a distintos escenarios sin riesgo para los equipos físicos.

    \item \textbf{Microsoft Excel:} Herramienta de apoyo para la organización de datos experimentales, análisis estadístico básico, procesamiento de resultados de simulaciones y generación de tablas y gráficas para la documentación del proyecto.

    \item \textbf{Módulos físicos de puente medio:} Proporcionados por la universidad colaboradora. Se utilizarán dos módulos para su conexión en paralelo bajo configuración interleaved, permitiendo implementar la topología estudiada de forma física y evaluar su comportamiento real.

    \item \textbf{Sistema de almacenamiento de energía (baterías y supercondensador):} Forman parte de la microrred DC existente y permitirán realizar pruebas prácticas de carga y descarga utilizando el convertidor implementado.

    \item \textbf{Fuente bidireccional DC/DC:} Fuente programable que permitirá inyectar o absorber potencia desde el sistema, simulando diferentes condiciones de operación y facilitando pruebas controladas del convertidor.

    \item \textbf{Instrumentación de laboratorio (osciloscopios, multímetros, fuente de alimentación, cargas electrónicas, etc.):} Utilizada para la medición de variables eléctricas clave (tensión, corriente, rizado), validación de señales de control y verificación del funcionamiento físico del convertidor implementado.
\end{itemize}

